\newpage
\pagestyle{plain}
\chapter{Туршилт, үр дүн ба хэлэлцүүлэг}

Энэхүү бүлэгт аппын туршилтын арга зүй, туршилтын үр дүн, мөн хэлэлцүүлгийг дэлгэрэнгүй тайлбарлана.

%% ============================================================
\section{Туршилтын арга зүй}
%% ============================================================

Аппын чанарыг баталгаажуулахын тулд дараах 3 түвшний туршилтын аргыг хэрэглэсэн.

\begin{table}[htbp]
\centering
\begin{tabular}{|p{0.8cm}|p{3.5cm}|p{4cm}|p{5cm}|}
\hline
\textbf{№} & \textbf{Туршилтын төрөл} & \textbf{Хэрэгсэл} & \textbf{Зорилго} \\
\hline
1 & Нэгж туршилт (Unit Test) & Flutter Test Framework & Загвар, үйлчилгээний логик шалгах \\
\hline
2 & Интеграцийн туршилт (Integration Test) & Flutter Integration Test & Дэлгэц хоорондын ажиллагаа шалгах \\
\hline
3 & Ашиглалтын туршилт (Usability Test) & Хэрэглэгчийн санал асуулга & UI/UX хэрэглэгчийн туршлага үнэлэх \\
\hline
\end{tabular}
\caption{Туршилтын арга зүйн хүрээ}
\label{tab:test_methods}
\end{table}

%% ============================================================
\section{Нэгж туршилт (Unit Testing)}
%% ============================================================

Нэгж туршилт нь аппын загвар классууд болон үйлчилгээний давхаргын логикийг бие даасан байдлаар шалгах зорилготой. Дараах туршилтуудыг гүйцэтгэсэн.

\subsection{Загвар классуудын туршилт}

Загвар класс бүрийн \texttt{fromMap()}, \texttt{toMap()}, \texttt{copyWith()} функцуудыг шалгасан.

\begin{table}[htbp]
\centering
\begin{tabular}{|p{3cm}|p{5cm}|p{2cm}|p{3.5cm}|}
\hline
\textbf{Загвар} & \textbf{Туршилт} & \textbf{Үр дүн} & \textbf{Тайлбар} \\
\hline
Person & fromMap() хөрвүүлэлт & Амжилттай & JSON $\rightarrow$ Person объект \\
\hline
Person & toMap() хөрвүүлэлт & Амжилттай & Person $\rightarrow$ Map\\
\hline
Person & copyWith() хуулбар & Амжилттай & Шинэ нэртэй хуулбар \\
\hline
Event & fromMap() хөрвүүлэлт & Амжилттай & JSON $\rightarrow$ Event объект \\
\hline
Quiz & fromMap() хөрвүүлэлт & Амжилттай & JSON $\rightarrow$ Quiz объект \\
\hline
MapData & fromMap() хөрвүүлэлт & Амжилттай & JSON $\rightarrow$ MapData объект \\
\hline
Media & fromMap() хөрвүүлэлт & Амжилттай & JSON $\rightarrow$ Media объект \\
\hline
\end{tabular}
\caption{Загвар классуудын нэгж туршилтын үр дүн}
\label{tab:model_tests}
\end{table}

\subsection{DataService-ийн туршилт}

\texttt{DataService}-ийн функцуудыг туршиж, дараах үр дүнг авсан.

\begin{table}[htbp]
\centering
\begin{tabular}{|p{4.5cm}|p{5cm}|p{2cm}|}
\hline
\textbf{Функц} & \textbf{Туршилтын нөхцөл} & \textbf{Үр дүн} \\
\hline
\texttt{loadAll()} & 5 JSON файл зэрэг ачаалах & Амжилттай \\
\hline
\texttt{searchPersons()} & ``Чингис'' гэж хайх & 1 үр дүн \\
\hline
\texttt{getPersonById()} & ID = 1 хайх & Чингис хаан \\
\hline
\texttt{getEventsForPerson()} & personId = 1 шүүх & 3+ үйл явдал \\
\hline
Давтан ачаалал & \texttt{loadAll()} 2 удаа дуудах & Нэг л удаа ачаална \\
\hline
\end{tabular}
\caption{DataService-ийн туршилтын үр дүн}
\label{tab:dataservice_tests}
\end{table}

%% ============================================================
\section{Функционал туршилт}
%% ============================================================

Аппын бүх үндсэн функцуудыг (FR-01 -- FR-07) гар аргаар болон автомат туршилтаар шалгасан.

\begin{table}[htbp]
\centering
\begin{tabular}{|p{1.2cm}|p{4cm}|p{5cm}|p{2cm}|}
\hline
\textbf{Код} & \textbf{Функц} & \textbf{Тест нөхцөл} & \textbf{Үр дүн} \\
\hline
FR-01 & Түүхэн хүмүүсийн намтар & Хүний жагсаалт, дэлгэрэнгүй мэдээлэл харах & Амжилттай \\
\hline
FR-02 & Газрын зураг & 3D бөмбөрцөг, 2D зураг, маркер дарах, zoom & Амжилттай \\
\hline
FR-03 & Цагийн хэлхээс & Үйл явдлын огноогоор эрэмбэлэлт, хүн рүү шилжих & Амжилттай \\
\hline
FR-04 & Квиз & Хариулт сонгох, оноо тооцох, үр дүн харах & Амжилттай \\
\hline
FR-05 & Мультимедиа & Зураг харуулах, карусель ажиллагаа & Амжилттай \\
\hline
FR-06 & Гэр бүлийн мод & 3 үеийн мод, zoom/pan, хүн дарах & Амжилттай \\
\hline
FR-07 & Соёлын мэдээлэл & Жагсаалт, дэлгэрэнгүй дэлгэц & Амжилттай \\
\hline
\end{tabular}
\caption{Функционал шаардлагын туршилтын үр дүн}
\label{tab:functional_tests}
\end{table}

%% ============================================================
\section{Гүйцэтгэлийн туршилт (Performance Testing)}
%% ============================================================

Аппын гүйцэтгэлийн үзүүлэлтүүдийг хэмжсэн.

\begin{table}[htbp]
\centering
\begin{tabular}{|p{5cm}|p{3cm}|p{3cm}|p{2.5cm}|}
\hline
\textbf{Үзүүлэлт} & \textbf{Хэмжилт} & \textbf{Зорилтот утга} & \textbf{Үнэлгээ} \\
\hline
Эхлэх хугацаа (Cold start) & $<$ 3 секунд & $<$ 5 секунд & Хангалттай \\
\hline
JSON ачаалах хугацаа & $<$ 500 мс & $<$ 1 секунд & Хангалттай \\
\hline
Дэлгэц хоорондын шилжилт & $<$ 300 мс & $<$ 500 мс & Хангалттай \\
\hline
3D бөмбөрцөг ачаалалт & $<$ 2 секунд & $<$ 3 секунд & Хангалттай \\
\hline
Санах ойн хэрэглээ & $\sim$120 МБ & $<$ 200 МБ & Хангалттай \\
\hline
Фреймийн хурд (FPS) & 55--60 fps & 60 fps & Хангалттай \\
\hline
\end{tabular}
\caption{Гүйцэтгэлийн туршилтын үр дүн}
\label{tab:performance_tests}
\end{table}

%% ============================================================
\section{Ашиглалтын туршилт (Usability Testing)}
%% ============================================================

Ашиглалтын туршилтыг 15 хэрэглэгчийн оролцоотойгоор гүйцэтгэсэн. Хэрэглэгчид нь 15--25 насны оюутан, сурагчид байсан.

\subsection{Туршилтын арга}

Хэрэглэгч бүрт дараах 5 даалгавар өгсөн:
\begin{enumerate}
	\item Чингис хааны намтрыг олж унших
	\item Газрын зургаас Самаркандын мэдээллийг олох
	\item Цагийн хэлхээсээс 1258 оны үйл явдлыг олох
	\item Квизэд оролцож, оноо авах
	\item Соёлын мэдээлэл хэсгээс нүүдлийн амьдралыг унших
\end{enumerate}

\subsection{Үнэлгээний үр дүн}

\begin{table}[htbp]
\centering
\begin{tabular}{|p{5cm}|p{2.5cm}|p{6cm}|}
\hline
\textbf{Шалгуур} & \textbf{Дундаж оноо (5-аас)} & \textbf{Тайлбар} \\
\hline
Хэрэглэхэд хялбар байдал & 4.3 & Ихэнх функцийг хялархан ашигласан \\
\hline
Дизайны нийцтэй байдал & 4.5 & Түүхийн уур амьсгал сайн илэрхийлсэн \\
\hline
Контентын чанар & 4.1 & Мэдээлэл үнэн зөв, товч тодорхой \\
\hline
Газрын зургийн интерактив байдал & 4.6 & 3D бөмбөрцөг маш сонирхолтой \\
\hline
Квизийн сонирхол татах байдал & 4.4 & Оноо, үр дүнгийн систем сайн \\
\hline
Нийт сэтгэл ханамж & 4.4 & Ерөнхийд нь сайн үнэлсэн \\
\hline
\end{tabular}
\caption{Ашиглалтын туршилтын үнэлгээний үр дүн}
\label{tab:usability_results}
\end{table}

\subsection{Хэрэглэгчийн санал хүсэлт}

Туршилтын явцад хэрэглэгчдээс дараах санал хүсэлтүүд гарсан:

\textbf{Эерэг санал:}
\begin{itemize}
	\item 3D газрын зураг маш сонирхолтой, түүхийг илүү ойлгомжтой болгосон
	\item Дизайн нь эртний бамбай цаасны загвартай, түүхийн уур амьсгалтай
	\item Квиз нь мэдлэгийг шалгахад тохиромжтой, оноо цуглуулах систем сонирхолтой
	\item Гэр бүлийн мод нь хаадын хоорондын хамаарлыг ойлгоход тусалсан
\end{itemize}

\textbf{Сайжруулах санал:}
\begin{itemize}
	\item Аудио тайлбар нэмвэл илүү сонирхолтой болно
	\item Квизийн асуултыг олшруулах, түвшин ялгах
	\item Англи хэл дээрх орчуулга нэмэх
	\item Favorite/Bookmark функц нэмэх
\end{itemize}

%% ============================================================
\section{Хэлэлцүүлэг}
%% ============================================================

\subsection{Зорилгын хэрэгжилт}

Төслийн эхэнд тавьсан зорилтуудын хэрэгжилтийг дүгнэвэл:

\begin{table}[htbp]
\centering
\begin{tabular}{|p{0.8cm}|p{5.5cm}|p{2cm}|p{5cm}|}
\hline
\textbf{№} & \textbf{Зорилт} & \textbf{Төлөв} & \textbf{Тайлбар} \\
\hline
1 & Түүхийн эх сурвалж судлах & Гүйцэтгэсэн & 10+ эх сурвалж ашигласан \\
\hline
2 & Шаардлага, дизайн тодорхойлох & Гүйцэтгэсэн & SRS, UI/UX, ERD бэлэн \\
\hline
3 & Flutter-ээр апп бүтээх & Гүйцэтгэсэн & 7 дэлгэц, 7 компонент \\
\hline
4 & Өгөгдөл хадгалах, sync бэлтгэх & Хэсэгчлэн & JSON оффлайн бэлэн, Firebase бэлтгэсэн \\
\hline
5 & Туршилт, сайжруулалт & Гүйцэтгэсэн & Unit, функционал, usability тест \\
\hline
6 & Үр дүн дүгнэх & Гүйцэтгэсэн & Энэхүү бүлэгт тусгасан \\
\hline
\end{tabular}
\caption{Зорилтуудын хэрэгжилтийн үнэлгээ}
\label{tab:goals_eval}
\end{table}

\subsection{Давуу ба сул талууд}

\textbf{Давуу талууд:}
\begin{itemize}
	\item Кросс-платформ хөгжүүлэлт — Android ба iOS-д нэг кодоор ажилладаг
	\item Оффлайн дэмжлэг — Сүлжээгүй орчинд бүрэн ажилладаг
	\item 3D интерактив газрын зураг — Ижил төрлийн аппуудаас ялгарах онцлог
	\item Гэр бүлийн модны визуалчлал — Хаадын хоорондын хамаарлыг ойлгомжтой харуулна
	\item Модульчлагдсан архитектур — Цаашид өргөтгөх, засварлахад хялбар
\end{itemize}

\textbf{Сул талууд:}
\begin{itemize}
	\item Онлайн синхрончлол бүрэн хэрэгжүүлэгдээгүй (Firebase бэлтгэсэн ч одоогоор оффлайн)
	\item Аудио, видео контент хязгаарлагдмал
	\item SQLite өгөгдлийн сан ашиглаагүй (JSON файлд суурилсан)
	\item iOS дээр туршилт хийгдээгүй (Android дээр голчлон туршсан)
\end{itemize}

\subsection{Ижил төрлийн системүүдтэй харьцуулалт}

\begin{table}[htbp]
\centering
\begin{tabular}{|p{3cm}|p{2.5cm}|p{2.5cm}|p{2.5cm}|p{2.5cm}|}
\hline
\textbf{Шалгуур} & \textbf{Энэ апп} & \textbf{Duolingo} & \textbf{Khan Academy} & \textbf{Монгол апп-ууд} \\
\hline
Интерактив & Өндөр & Өндөр & Дунд & Бага \\
\hline
3D газрын зураг & Тийм & Үгүй & Үгүй & Үгүй \\
\hline
Оффлайн & Бүрэн & Хэсэгчлэн & Хэсэгчлэн & Ихэвчлэн \\
\hline
Монгол хэл & Бүрэн & Үгүй & Үгүй & Бүрэн \\
\hline
Квиз систем & Тийм & Тийм & Тийм & Цөөн \\
\hline
\end{tabular}
\caption{Ижил төрлийн системүүдтэй харьцуулалт}
\label{tab:comparison}
\end{table}
